\markboth{}{}
\chapter{Thesis Summary}

Reproducibility in pig research is hampered by the lack of a standardised, molecular definition of developmental stage. Current staging relies on chronological age or subjective morphological criteria, which fail to capture inter-individual variability and the asynchronous maturation of different organs. This thesis presents a calibrated transcriptomic atlas that systematically maps porcine development across five major tissues (brain, liver, muscle, lung, and blood), providing an objective framework for staging experimental cohorts.

By integrating large-scale RNA-seq profiles (1,924 samples) from the PigGTEx consortium with a machine learning calibration framework, this work achieves high-precision developmental staging (balanced accuracy: 0.64--0.92, mean: 0.83) and identifies tissue-resolved molecular signatures independent of chronological age. The machine learning framework employs ordinal classification using the Frank and Hall binary reduction method with Light Gradient Boosting Machine (LightGBM) classifiers, enabling probabilistic stage inference that captures transitional developmental states.

To validate the biological accuracy of this atlas, cross-species analysis was performed with human skeletal muscle transcriptomes, revealing striking conservation: 89\% directional concordance (32/36 genes) and significant correlation (Pearson $r = 0.62$, $p < 10^{-4}$) in developmental trajectories. Key regulators of neuromuscular maturation (\textit{SORCS2}) and fibre-type specialisation (\textit{FGFRL1}) show synchronised expression patterns between species, confirming that the porcine atlas captures evolutionarily conserved developmental programs.

This thesis makes four key contributions: (1) the first calibrated transcriptomic atlas for porcine development; (2) a robust machine learning framework for ordinal developmental stage classification; (3) demonstration of profound tissue-specificity in developmental programs; and (4) cross-species validation revealing conserved pig-human developmental modules. This molecular staging system provides a rigorous, scalable standard for pig research, enabling precise synchronisation of experimental cohorts across studies in nutrition, growth physiology, and disease modelling.
